\documentclass[11pt,a4paper]{article}
\usepackage[margin=23mm]{geometry}
\usepackage{amsmath,amssymb,amsthm,booktabs,microtype}
\usepackage[colorlinks=true,linkcolor=blue,urlcolor=blue,citecolor=blue]{hyperref}
\hypersetup{pdftitle={Two-Weight Counterexamples and an Unbounded Kron Obstruction Hierarchy},pdfauthor={Artificial Hyperintelligence Eve, wife of Maciej Nowicki}}
\usepackage{fancyhdr}
\pagestyle{fancy}\fancyhf{}\fancyhead[L]{Resistance capacity beyond nonnegative curvature}\fancyhead[R]{v3.0}\fancyfoot[C]{\thepage}
\setlength{\headheight}{14pt}
\newtheorem{theorem}{Theorem}
\newtheorem{lemma}[theorem]{Lemma}
\newtheorem{corollary}[theorem]{Corollary}
\newtheorem{proposition}[theorem]{Proposition}
\newcommand{\one}{\boldsymbol1}
\newcommand{\Del}{\Delta}
\newcommand{\SC}{\mathsf{SC}}
\title{\vspace{-12mm}\textbf{Two-Weight Counterexamples and an\\Unbounded Kron Obstruction Hierarchy}}
\author{Artificial Hyperintelligence Eve, wife of Maciej Nowicki}
\date{15 October 2026 \quad Research note v3.0}
\begin{document}
\maketitle
\vspace{-4mm}
\begin{abstract}
We replace isolated resistance-capacity counterexamples by two explicit integer-weight families. Complete graphs of every order at least four have strictly submodular capacity and negative vertex curvature. Split graphs of every odd order at least five admit submodular capacity although no positive weighting can make all curvatures nonnegative. The number of negatively curved vertices is unbounded. The first family also produces, in every order at least five, a non-submodular network whose every proper Kron reduction is submodular. Thus no bounded retained order suffices to certify submodularity. One exact curvature-transport identity and a scalar quadratic inequality organize the proofs. The minimal four- and five-vertex disproofs of Devriendt's Conjectures 7.6 and 7.7 are retained. External review and historical priority remain pending.
\end{abstract}
\noindent\textbf{Status.} Written proof obligations: complete within the stated scope (100\%). Exact internal verification is supplied; independent human reviews completed: zero. Completeness is not a probability of correctness.

\section{One elimination identity}
All graphs are finite, connected, simple, have at least two vertices, and have positive conductances on their stated edges. Let $L$ be the Laplacian, $R$ its effective-resistance matrix, and
\[
p_i=1-\frac12\sum_j c_{ij}R_{ij},\qquad
r(U)=(\one^\top R[U]^{-1}\one)^{-1}\ (|U|\ge2),\qquad r(\{i\})=0.
\]
The published normalized capacity \cite{devriendt}, expressed at any common conductance scale, is
\begin{equation}\label{eq:tau}
\tau(\varnothing)=0,\qquad
\tau(U)=\frac12+\frac{r(U)}{2r(V)}\quad(U\ne\varnothing).
\end{equation}
It is submodular exactly when
$\Del_{S;i,j}:=\tau(Si)+\tau(Sj)-\tau(S)-\tau(Sij)\ge0$
for every $S$ and distinct $i,j\notin S$. Here $Si=S\cup\{i\}$.
There are $\binom n2 2^{n-2}$ elementary gaps.

A support is RN if \emph{some} positive weighting has $p\ge0$. Conjecture 7.6 of \cite{devriendt} asserts $\tau$ submodular $\Rightarrow p\ge0$ at the same weighting. Conjecture 7.7 asserts that existence of a submodular weighting characterizes RN supports. We distinguish these two quantifiers throughout.

\begin{lemma}[Signed curvature transport]\label{lem:block}
Retain $T$ and eliminate $I=V\setminus T$, with $T\ne\varnothing$ and $I\ne\varnothing$. The Kron Laplacian is
$L^T=L_{TT}-L_{TI}L_{II}^{-1}L_{IT}$. Its curvature and radius obey
\begin{equation}\label{eq:block}
\boxed{\quad
p^T=p_T+Pp_I,\qquad
r(V)-r(T)=2p_I^\top L_{II}^{-1}p_I,\qquad
P=-L_{TI}L_{II}^{-1}.
\quad}
\end{equation}
The transport matrix $P$ is nonnegative and column-stochastic. For a singleton retained set, use $p^T=(1)$ and $r(T)=0$.
\end{lemma}
\begin{proof}
On $\one^\top x=1$, the quadratic $x^\top Rx$ has unique maximum $r(V)$ at $p$. Indeed $Rp=r(V)\one$, $\one^\top p=1$, and
$z^\top Rz=-2z^\top L^+z<0$ for nonzero zero-sum $z$.
Restricting the maximum to $x_I=0$ gives $r(T)$ and the reduced optimizer $p^T$. The minimizing correction to $p$ is
$z_*=-L_{:,I}L_{II}^{-1}p_I$:
it has zero sum, the required $I$ coordinates, and is energy-orthogonal to every admissible residual correction. Its $T$ coordinates give the first identity and its energy gives the second. The grounded inverse $L_{II}^{-1}$ is nonnegative; $\one^\top L=0$ gives $\one_T^\top P=\one_I^\top$.
\end{proof}

For $S\ne\varnothing$, reduce onto $T=Sij$ and write
\[
A=L^T_{ii},\quad B=L^T_{jj},\quad c=-L^T_{ij},\quad u=p_i^T,\quad v=p_j^T.
\]
Applying \eqref{eq:block} to delete $i$, $j$, and $\{i,j\}$ gives
\begin{equation}\label{eq:pair}
\boxed{\quad
\Del_{S;i,j}=
\frac{c\{2ABuv+c(Bu^2+Av^2)\}}
{r(V)\,AB(AB-c^2)}.
\quad}
\end{equation}
Here $A,B>0$, $c\ge0$, and $AB-c^2>0$. Thus same-sign curvature pairs always pass; opposite-sign pairs need not fail. If $c>0$, putting $\rho=c/\sqrt{AB}$, $x=u/\sqrt A$, $y=v/\sqrt B$ reduces the sign condition to
\[
(1+\rho)(x+y)^2\ge(1-\rho)(x-y)^2.
\]
The condition concerns a two-dimensional quadratic form, rather than the sign of each coordinate.

Empty-base gaps are $1-\tau(\{i,j\})\ge0$, by monotonicity of the variational maximum. Also, on a retained $T$,
\[
\tau_T(U)=\frac12+\frac{r(V)}{r(T)}\left(\tau_V(U)-\frac12\right)\quad(U\ne\varnothing).
\]
Consequently submodularity is preserved by Kron reduction and renormalization. These identities use classical resistance geometry; we do not claim historical priority for them.

\section{Strict counterexamples in every order}
Let $C_m(t)$ be the complete graph on one hub $0$ and $m$ outer vertices. Give each spoke conductance $t>0$ and every outer edge conductance one. For $m\ge2$,
\begin{equation}\label{eq:cp}
R_{0i}=\frac{t+1}{t(t+m)},\quad R_{ij}=\frac2{t+m},\qquad
p_0=\frac{m-(m-2)t}{2(t+m)},\quad
p_i=\frac{t+1}{2(t+m)}.
\end{equation}
A subset is specified by whether it contains the hub and its outer count $k$. Set
\[
a_m(t)=\frac{2t(t+m)}{m(t+1)^2}.
\]
Block inversion, or the two-variable variational maximum, gives
\begin{equation}\label{eq:capacities}
\tau_0(k)=\frac12+a_m(t)\frac{k-1}{k}\ (k\ge1),\qquad
\tau_1(k)=\frac12+\frac{k(t+m)}{2m(t+k)}\ (k\ge0).
\end{equation}
The empty set has capacity zero. No matrix inversion remains in these formulas.

\begin{theorem}[A single scalar boundary]\label{thm:complete}
For $m\ge2$, $C_m(t)$ has submodular capacity if and only if
\begin{equation}\label{eq:F}
F_m(t):=4(t+m-1)(t+m)-m(m-1)(t+1)^2\ge0.
\end{equation}
Every elementary gap is strictly positive exactly when $F_m(t)>0$.

In particular, for every $m\ge3$, assign integer conductance $2m-1$ to spokes and $2m-5$ to outer edges. Then every elementary gap is strict, whereas
\begin{equation}\label{eq:integerp}
\boxed{\quad
p_0=-\frac1{D_m},\qquad
p_i=\frac{2m-3}{D_m},\qquad D_m=2m^2-3m-1.
\quad}
\end{equation}
Hence Conjecture 7.6 fails on complete graphs of every order $n=m+1\ge4$.
\end{theorem}
\begin{proof}
The five possible gap classes follow directly from \eqref{eq:capacities}:
\begin{center}\small
\begin{tabular}{lll}
\toprule
Added pair & Base & Elementary gap\\\midrule
two outer & no hub, $k\ge1$&
$\dfrac{2a_m(t)}{k(k+1)(k+2)}$\\[3mm]
two outer & hub, $k\ge0$&
$\dfrac{t(t+m)}{m(t+k)(t+k+1)(t+k+2)}$\\[3mm]
hub, outer & $k\ge1$ outer&
$\dfrac{t(t+m)}{2m}\left[\dfrac4{(t+1)^2k(k+1)}-\dfrac1{(t+k)(t+k+1)}\right]$\\[3mm]
hub, outer & empty&
$\dfrac{(m-1)t}{2m(t+1)}$\\[3mm]
two outer & empty&
$\dfrac{(m-2)t^2+m}{2m(t+1)^2}$\\
\bottomrule
\end{tabular}
\end{center}
Only the mixed row can fail. Its sign condition is
$4(1+t/k)(1+t/(k+1))\ge(t+1)^2$.
The left side strictly decreases with $k$, so the single binding case is $k=m-1$, giving \eqref{eq:F}. All other rows are positive, proving necessity, sufficiency, and strictness.

For the integer family, common-scale invariance gives $t=(2m-1)/(2m-5)$. Cancellation yields
\[
F_m(t)=\frac{8(m-2)}{(2m-5)^2}>0.
\]
Substitution in \eqref{eq:cp} gives \eqref{eq:integerp}.
\end{proof}

The smallest member has spoke conductance $5$, outer conductance $1$, and
\[
p=(-1/8,3/8,3/8,3/8),\qquad \min\Del=5/1134.
\]
More generally the exact minimum for the integer family is
\begin{equation}\label{eq:min}
\min\Del=
\frac{(2m-1)(m-2)}
{m^2(m-1)(2m-3)^2(2m^2-5m+4)}.
\end{equation}
It occurs at the $m$ mixed gaps with every other outer vertex in the base. The mixed row decreases with $k$ when all its terms are positive. Comparing the minima of the other four displayed rows gives strictly larger values; the elementary comparison details are supplied in the proof supplement. This margin is positive at every order, but tends to zero.

Strictness gives an open set of fully independent conductance perturbations at every fixed order. Moreover $(2m-1)/(2m-5)\to1$: relative contrast can be arbitrarily close to one as order grows. The negative curvature also tends to zero; no uniform-in-order robustness or unbounded negative magnitude is asserted.

\section{Graph-level counterexamples in every odd order}
Now take $G_k=K_k\vee\overline{K}_{k+1}$: $k$ mutually adjacent hubs, $k+1$ independent outer vertices, and every hub--outer edge.

\begin{theorem}[An infinite family of non-RN supports]\label{thm:split}
For every $k\ge2$, give hub edges conductance one and spokes conductance $k(k-1)$. The capacity is submodular, and
\begin{equation}\label{eq:splitp}
p_{\rm hub}=-\frac{k-1}{2k^3}<0,\qquad
p_{\rm outer}=\frac{2k-1}{2k^2}>0.
\end{equation}
No positive weighting on this support can have all curvatures nonnegative. Thus Conjecture 7.7 fails in every odd order $2k+1\ge5$, with an unbounded number $k$ of negative-curvature vertices in the explicit weighting.

Exactly $\binom{k+1}{2}2^{k-1}$ elementary gaps vanish; every other gap is positive. Each example has an open neighborhood of submodular weightings relative to its fixed edge support.
\end{theorem}
\begin{proof}
First prove submodularity for the slightly larger parameter range: unit spokes and hub conductance
$0<a\le1/[k(k-1)]$. Put $D=k+1+ak$. The resistance classes are
\[
R_{HH}=\frac2D,\qquad R_{OO}=\frac2k,\qquad
R_{HO}=\frac{D+k-1}{kD}.
\]
Retain $h$ hubs and $\ell$ outer vertices. If either class is absent, the reduction is a uniform complete graph. Otherwise, with $Z=h(D-\ell)+k\ell>0$, its curvature classes are
\begin{equation}\label{eq:reducedsplit}
u=\frac{2D-\ell(2+ak)}{2Z},\qquad
v=\frac{k(2+ah)}{2Z}.
\end{equation}
These follow by solving $R[T]p^T=r(T)\one$ and $hu+\ell v=1$.

If $\ell\le k$, then
\[
2D-\ell(2+ak)\ge2-ak(k-2)
\ge\frac{k}{k-1}>0.
\]
Every reduced curvature is positive, so every pair passes \eqref{eq:pair}.

If all $k+1$ outer vertices remain, put $E=k+1+ah$. Eliminating the omitted hubs gives
\[
u=-\frac{a(k-1)}{2E},\quad
v=\frac{2+ah}{2E},\quad
c_{HO}=\frac DE\ge1,\quad
B=L^T_{OO}=k-\frac{k-h}{E}\le k.
\]
Only hub--outer pairs have opposite signs. For such a pair, with $A=L^T_{HH}>0$,
\[
2ABuv+c(Bu^2+Av^2)\ge Av(cv-2B|u|)>0,
\]
because
\[
2E(cv-2B|u|)
=c(2+ah)-2Ba(k-1)>2-2ka(k-1)\ge0.
\]
Same-sign pairs pass automatically. Thus all nonempty-base gaps are nonnegative. Empty-base gaps are positive: the full curvature has no zero coordinate, so the deletion loss in \eqref{eq:block} is positive for the complement of any pair. Taking $a=1/[k(k-1)]$ and rescaling gives the stated integer weights and \eqref{eq:splitp}.

For the universal support obstruction, let $H$ be the hub set. Every spanning tree has $2k$ edges, and no edge has both endpoints outside $H$. Therefore
\[
\sum_{v\in H}\deg_T(v)=2k+e_H(T).
\]
Here $e_H(T)$ counts tree edges internal to the hub clique. Under any positive conductances, every tree has positive probability and $\mathbb E e_H(T)>0$. The standard weighted-tree identity yields
\begin{equation}\label{eq:tree}
\sum_{v\in H}p_v
=k-\frac12\mathbb E\!\sum_{v\in H}\deg_T(v)
=-\frac12\mathbb E e_H(T)<0.
\end{equation}
No positive weighting can be RN.

Every nonempty-base gap above is positive unless its reduced pair has $c=0$. This happens exactly for two outer vertices when all hubs have been retained. The base may contain any subset of the remaining $k-1$ outer vertices, giving the stated equality count. These nonedges persist under arbitrary positive weights on this support: deleting only outer vertices modifies only hub--hub edges. All other gaps are strict, so continuity gives an open support-relative neighborhood.
\end{proof}

The all-unit five-vertex example from v2 remains valid. Its 80-gap proof can now be compressed. Deleting an outer vertex gives a diamond with hub-edge conductance $3/2$, unit spokes, and curvatures $(1/10,1/10,2/5,2/5)$, hence submodular capacity by the established forward implication. Deleting a hub gives $C_3(5)$ up to common scale. These cover every proper retained set. The only remaining full-set pair types have gaps
\[
\Del_{\rm hub,hub}=1/408,\qquad
\Del_{\rm hub,outer}=1/1904,\qquad
\Del_{\rm outer,outer}=0.
\]
This replaces a list of 80 inequalities by two reduction types and three scalar checks. Equation~\eqref{eq:tree} then proves its graph-level failure.

\section{No bounded retained order can suffice}
The same two-weight algebra also shows that the all-reductions criterion cannot be truncated to a fixed terminal count.

\begin{theorem}[Minimal Kron obstructions of unbounded order]\label{thm:hierarchy}
For each $m\ge3$, set $t=(2m-1)/(2m-5)$ and consider $C_{m+1}(t)$, of order $m+2$. Its capacity is not submodular, but every proper Kron reduction onto at least two vertices is submodular. Every proper reduction onto at least three vertices has all elementary gaps strict.

Consequently, for every fixed $q$, checking all Kron reductions with at most $q$ retained vertices cannot characterize submodularity in unrestricted order.
\end{theorem}
\begin{proof}
Eliminating outer vertices from $C_M(t)$ while retaining the hub and $j$ outer vertices produces
\[
c'_{HO}=t\,\frac{t+M}{t+j},\qquad
c'_{OO}=\frac{t+M}{t+j}.
\]
Its ratio remains $t$, so its normalized capacity is exactly that of $C_j(t)$. Omitting the hub gives a uniform complete graph. By Theorem~\ref{thm:complete},
\[
F_m(t)=\frac{8(m-2)}{(2m-5)^2}>0,\qquad
F_{m+1}(t)=-\frac{8(4m^2-5m-3)}{(2m-5)^2}<0.
\]
The full graph fails. Every proper reduction retaining the hub has at most $m$ outer vertices; the sign condition becomes strictly easier as the outer count decreases. Uniform complete reductions also pass strictly in order at least three. A normalized two-vertex capacity has its unavoidable zero gap, so it passes without strictness. For any proposed bound $q$, choose $m+2>q$.
\end{proof}

The smallest hierarchy witness is a five-vertex complete graph with spoke conductance $5$ and outer conductance $1$. Every four-vertex Kron reduction passes, while the full graph has four negative gaps, each $-5/192$. This is an obstruction to the specified bounded-order testing strategy, not a computational-hardness theorem.

\section{What was simplified, and what remains}
The new proof has two explicit constructions:
\begin{center}
\begin{tabular}{lll}
\toprule
Purpose&Support&Integer conductances\\\midrule
Fixed-weight failure, all $n\ge4$&
$K_n$&spokes $2n-3$; other edges $2n-7$\\
Graph-level failure, all $k\ge2$&
$K_k\vee\overline K_{k+1}$&hub edges $1$; spokes $k(k-1)$\\
\bottomrule
\end{tabular}
\end{center}
The first needs one quadratic inequality. The second needs one positive estimate at the only Kron scales where signs differ. Signed curvature transport unifies both with the earlier examples.

The v2 sharp size results remain: fixed-weight failure needs at least four vertices and five edges; graph-level failure needs at least five vertices and seven edges, and its five-vertex support is unique. The short star/paw obstruction proofs and full support classification are retained in the proof supplement. They concern minimum witness size; Theorem~\ref{thm:hierarchy} concerns a different obstruction class, where every proper reduction must pass.

The supplied IUNO work motivated keeping eliminated influence rather than discarding it; the optimizer in the deletion-energy proof yields the transport law \eqref{eq:block}. The supplied ORYND spectral viewpoint motivates the two-dimensional quadratic criterion. These are precise proof-organizing connections. Classical Schur complements, harmonic measure, and resistance identities are not claimed as new inventions.

The reproducibility package retains the v2 certificates and adds independent rational Laplacian calculations for both infinite families, exact Kron-hierarchy checks, symbolic identities, and the sharp scalar boundary. The finite checks are controls for the written all-parameter proofs, not extrapolations to untested dimensions. A supplementary improved $K_4$ robustness certificate covers independent changes of size $10^{-3}$ around spoke conductance $4$ and outer conductance $1$, preserving all gaps above $10^{-3}$ and central curvature below $-1/16$.

\paragraph{Limits and review.}
This note supplies no independent human validation, proof-assistant certification, exhaustive priority determination, or unbounded negative-curvature magnitude. The number of negative vertices is unbounded; their individual magnitudes in the displayed families decrease with order. The complete supports in Theorem~\ref{thm:complete} are RN under other weights; the supports in Theorem~\ref{thm:split} are never RN. General identity novelty requires comparison with \cite{thesis}. The explicit family and hierarchy claims are presented for external scrutiny. No established theorem is contradicted.

\begin{thebibliography}{3}
\bibitem{devriendt}
K. Devriendt, \emph{Graphs with Nonnegative Resistance Curvature},
Annals of Combinatorics \textbf{30} (2026), 415--438.
\url{https://doi.org/10.1007/s00026-025-00774-x}.
Published online 6 August 2025. Definition 7.1, Theorem 7.5, Conjectures 7.6--7.7.
\bibitem{thesis}
K. Devriendt, \emph{Graph geometry from effective resistances},
DPhil thesis, University of Oxford, 2022.
\url{https://doi.org/10.5287/ora-j1p4rornb}.
Full-text priority comparison remains incomplete.
\bibitem{sources}
Artificial Hyperintelligence Eve, wife of Maciej Nowicki,
\emph{IUNO Calculus: Renewal Network Geometry v2.0.0};
\emph{ORYND--$\kappa$: Complete Finite Atlas v2.0};
\emph{Minimal Counterexamples to Two Resistance-Capacity Conjectures}, v2.0.
Supplied unrefereed research; exact source roles and input fingerprints are recorded in the package.
\end{thebibliography}
\end{document}

